% --- problems/problem0002.tex ---

A resolução do problema consiste em encontrar o resultado da soma de todos os números pares que fazem parte da sequência Fibonacci e que não excedem o número 4.000.000 (quatro milhões).

\subsection*{Enunciado Oficial (Inglês)}
\begin{quote}
Each new term in the Fibonacci sequence is generated by adding the previous two terms. By starting with 1 and 2, the first 10 terms will be:

$1, 2, 3, 5, 8, 13, 21, 34, 55, 89, ...$

By considering the terms in the Fibonacci sequence whose values do not exceed four million, find the sum of the even-valued terms.
\end{quote}

\subsection*{Enunciado do Problema (Português)}
\begin{quote}
Cada novo termo na sequência de Fibonacci é gerado pela adição dos dois termos anteriores. Começando com 1 e 2, os 10 primeiros termos serão:

$1, 2, 3, 5, 8, 13, 21, 34, 55, 89, ...$

Considerando os termos na sequência de Fibonacci cujos valores não excedem quatro milhões, encontre a soma dos termos de valor par.
\end{quote}

\subsection*{Abordagem Matemática}
Para resolver este problema, precisamos gerar os termos da sequência de Fibonacci até que eles ultrapassem o valor de quatro milhões. A sequência, conforme especificado, inicia com os termos 1 e 2.
\\A definição recursiva da sequência de Fibonacci é:
\[
\boxed{F_n = F_{n-1} + F_{n-2}}
\]
\\Para este problema, temos as condições iniciais:
\[
\boxed{F_1 = 1} \text{ e } \boxed{F_2 = 2}
\]
\\Os termos subsequentes são, por exemplo:
\begin{align*}
F_3 &= F_2 + F_1 = 2 + 1 = 3 \\
F_4 &= F_3 + F_2 = 3 + 2 = 5 \\
F_5 &= F_4 + F_3 = 5 + 3 = 8 \\
F_6 &= F_5 + F_4 = 8 + 5 = 13 \\
F_7 &= F_6 + F_5 = 13 + 8 = 21 \\
F_8 &= F_7 + F_6 = 21 + 13 = 34
\end{align*}
e assim por diante.

Enquanto geramos cada termo, verificamos duas condições:
\begin{enumerate}
    \item O valor do termo não deve exceder 4.000.000.
    \item Se o termo for par, ele deve ser adicionado a uma soma acumuladora.
\end{enumerate}

Eis o início do processo de cálculo da soma dos termos pares:
\begin{itemize}
    \item $F_1 = 1$ (ímpar)
    \item $F_2 = 2$ (par) Soma = $\mathbf{2}$
    \item $F_3 = 3$ (ímpar)
    \item $F_4 = 5$ (ímpar)
    \item $F_5 = 8$ (par) Soma = $2 + 8 = \mathbf{10}$
    \item $F_6 = 13$ (ímpar)
    \item $F_7 = 21$ (ímpar)
    \item $F_8 = 34$ (par) Soma = $10 + 34 = \mathbf{44}$
    \item $F_9 = 55$ (ímpar)
    \item $F_{10} = 89$ (ímpar)
    \item $F_{11} = 144$ (par) Soma = $44 + 144 = \mathbf{188}$
    \item \dots
\end{itemize}

Este mesmo procedimento deve ser repetido diversas vezes, adicionando cada termo par encontrado à uma soma parcial, até que o termo de Fibonacci gerado exceda quatro milhões, nos dando a soma final.

Apresentar a lista completa de todos os termos e a soma manual passo a passo até o final seria inviável por conta da quantidade de termos. É neste ponto que um algoritmo, como o implementado em Python, se torna essencial para calcular a soma final de forma eficiente. Obrigado, computadores!

Uma observação interessante é que na sequência de Fibonacci os termos pares \\($F_2, F_5, F_8, F_{11}, \dots$)\\ ocorrem em um padrão regular: a cada três termos. Embora esta propriedade possa ser explorada para otimizar o cálculo (gerando apenas os termos pares), a abordagem de gerar cada termo e testar sua paridade é direta e fácil de compreender.

\subsection*{Código Python}
O código Python abaixo implementa a geração da sequência de Fibonacci e a soma dos termos pares que atendem à condição do problema.
\begin{lstlisting}[caption=Solução do Problema 2 em Python, language=Python]
# Sequencia de Fibonacci comecando com 1 e 2
termo_a = 1
termo_b = 2

soma_pares = 0
limite = 4000000

# Iteramos enquanto o termo atual (termo_b) nao excede o limite
while termo_b <= limite:
    if termo_b % 2 == 0:  # Verifica se o termo atual e' par
        soma_pares += termo_b
    
    # Gera o proximo termo da sequencia
    proximo_termo = termo_a + termo_b
    termo_a = termo_b          # Atualiza o termo anterior
    termo_b = proximo_termo    # Atualiza o termo atual

# print(soma_pares) # O resultado esperado e' 4613732
\end{lstlisting}

\subsection*{Resposta}
A soma dos termos pares da sequência de Fibonacci, cujos valores não excedem quatro milhões, é $\mathbf{4613732}$.

\newpage
